The full-state comparison in Table~\ref{tab:decoupling_audit} changes a
supervision bundle, not the presence versus absence of training. The audited
legacy route applies weighted native reconstruction, derivatives, sparse FK,
contact, boundary, endpoint and residual penalties. Its q-only reconstruction
uses weighted Huber loss over all 34 native channels: root channels have weight
2.0, leg channels 1.5, waist channels 1.25, and the remaining arm channels 0.5.
Thus fine native articulation remains in the structural objective. That term
also includes native first-difference loss (0.5) and contact BCE (0.1), before
its outer coefficient of 0.5.

The frozen legacy trainer specifies full-reconstruction coefficients of 1.0
for native motion, 0.5 for velocity, 0.1 for acceleration, 0.5 for FK, 0.1 for
contact BCE, 0.2 for contact height, 0.1 for contact slide, 1.0 for VQ, and 0.2
for the amplitude guard. Boundary loss has outer weight 1.0 and internal
velocity/acceleration/jerk weights 1.0/0.5/0.25. Additional coefficients are
0.5 for the q-only endpoint-state loss, 0.25 for critical-state invariance,
and 0.01 for residual rate. Training corruption keeps the residual unchanged
with probability 0.55, removes it with probability 0.30, and adds channel-scaled
noise with probability 0.15. On this legacy path, corrupted reconstructions
can still be trained toward the full target.

In contrast, the proposed supervision uses only the four terms of
Eq.~\eqref{eq:codec_loss}. Full reconstruction uses normalized native motion,
boundary-aware velocity, all-link/sole geometry and geometry velocity.
The discrete and corrupted-residual reconstructions are supervised only on
the declared structure projector, excluding its exact detail complement.
Each projector loss equally averages its four normalized blocks. The outer
coefficients are 1.0 for full reconstruction, 1.0 for discrete structure,
0.25 for residual robustness, and 1.0 for VQ; the VQ commitment coefficient
is 0.25. The contact head and the separate legacy physical, endpoint,
invariance, amplitude, and residual-rate penalties are absent. Physical
diagnostics remain evaluation measures rather than hidden extra losses.

The source definitions and frozen experiment contract establish these
supervision differences. Exact checkpoint/launch-argument binding remains part
of the historical table provenance audit. This comparison supports the effect
of the supervision bundle; attributing it to any one removed penalty would
require an additional controlled ablation.
